\chapter{Conclusion}
\label{sec:conclusion}

This thesis started from a gap in how vision-language reward models are evaluated. The field reports strong offline metrics and selected reinforcement learning results, real-world runs or simulated runs after in-domain adaptation, and the most intuitive experiment, a released foundation reward model driving simulated RL as it stands, was missing. We built RoboRef for exactly that test, a reward model meant to perform well while a policy optimizes against it, through a re-curated corpus whose failures carry dense per-frame progress labels, an optional in-context demonstration alongside every query, and an asymmetric objective that makes false positives the expensive mistake. The work was organized around four short questions. Can a reward model trained mostly on successful demonstrations handle the failures a policy generates, or is failure data needed (RQ1)? Do in-context demonstrations improve its judgment (RQ2)? Are these models overconfident, and does an asymmetric loss make them harder to exploit (RQ3)? And do offline gains transfer to downstream reinforcement learning (RQ4)? We can now answer each in turn.

\textbf{RQ1}, whether failure supervision improves the discrimination of success from failure, has a clear answer in distribution. Changing nothing but the training data lifts rank accuracy from $0.611$ to $0.771$ and doubles the true-positive rate at a strict false-positive budget. The ablation of Section~\ref{sec:results-icl} answers the question exactly from the other side, since continuing to train the released model with our loss but on its own data moves rank accuracy only from $0.61$ to $0.64$, against the $0.85$ the same loss reaches with our corpus, so the gain belongs to the failure data and not to the objective wrapped around it. On Robomimic the standard-loss fine-tune, which differs from the released model in the training data alone, clears behavior cloning on a domain where the released model never moves at all, the data's contribution measured on-policy. That it cannot hold the gain without the asymmetric loss makes the division of labor plain. The data is the foundation the other two contributions stand on, necessary even where it is not sufficient. However, it should be noted that the best RL results come from the success head, which the dense per-frame labels do not directly supervise. Whether the reason behind this is the training loss of the progress head or our per-frame labeled dataset remains an open question.

\textbf{RQ2}, whether in-context demonstrations improve the model's judgment, has to be answered by the standard this thesis has insisted on. Offline the demonstration is the clearest yes we measured. It lifts rank accuracy to $0.896$, more than $27$ points above the released baseline, with the widest and still-growing margin on the real-world split, exactly where a textual task description is thinnest. We have argued throughout, however, that offline metrics are not where a reward model earns its keep, and on-policy the case does not hold. Across MetaWorld, LIBERO, and Robomimic the ICL variant is never clearly better than the plain asymmetric fine-tune, while the demonstration doubles the input the model reads on every query and with it the cost of deployment. By our own standard the answer is therefore negative. The demonstration sharpens judgment on a fixed test set but does not buy a better policy, and at twice the inference cost it is not justified in the loop as we deployed it. That the largest offline gain in the thesis does not survive the optimizing policy already previews the answer to RQ4.

\textbf{RQ3} asked whether these models are overconfident and whether an asymmetric loss makes them harder to exploit. Both halves held. The released baseline scores genuine failures next to genuine successes and, under an optimizing policy, pays out the most reward in the comparison while training a policy that never succeeds. The asymmetric loss is the single largest downstream factor we measured, the whole distance between a fine-tune that collapses under an optimizing policy and one that reaches $0.72$ on MetaWorld and ends above behavior cloning on a domain it never saw.

\textbf{RQ4}, whether offline gains transfer to downstream RL, gets the answer no, not by themselves, and this is the finding we consider most important beyond our own models. Offline metrics score a fixed distribution while a policy actively searches for the reward's mistakes, so a model can match the best fine-tunes on nearly every offline column and still be pulled apart on-policy. What transfers is conservatism. The failure mechanism, visible because simulation provides ground truth, is a policy farming false positives, collecting reward that rises while true success does not.

The main question, to what extent a fine-tuned vision-language model can serve as a usable and generalizable reward for robot RL, therefore gets a conditional but genuinely positive answer. With failure supervision and a conservative objective, one reward model trained policies across every simulated domain we tested, including one fully outside its training distribution, and delivered the strongest overall downstream record of any reward model in the comparison. 

However, deployment needs a calibrated threshold and one demonstration-length prior, results inherit the sensitivity of both the model and the RL algorithm beneath it, and everything we show lives in simulation, with seed counts bounded by compute. A real-world learning loop, online threshold calibration, and reward pretraining on newer backbones are where this work should go next. What the thesis establishes is the direction: a foundation reward model earns that name in the RL loop, not on a leaderboard, and the way to survive the loop is to make the reward hard to exploit.
